\documentclass{article}
\usepackage{mathptmx}
\usepackage{courier}
\usepackage[paper=a4paper]{geometry}
\usepackage{amsmath,amssymb,amsfonts,mathtools}
\usepackage{fontawesome5}
\usepackage{hyperref}
\usepackage{tabularx}
\usepackage{enumitem}
\usepackage{multicol}

\hypersetup{
	colorlinks=true
}

\newcommand{\code}[1]{\texttt{#1}}
\newcounter{problem}
\NewDocumentEnvironment{problem}{o}{
	\stepcounter{problem}
	\paragraph{Problem \theproblem}
	\IfValueT{#1}{ (#1 Points)}
}{}

\title{Multiple Linear Regression Practice Set}
\author{Bhaya A.G.\thanks{\faEnvelope[regular]\ Electronic mail: \href{mailto:angubh.io@gmail.com}{\code{ananya[dot]gb[at]nbu[dot]ac[dot]in}}}}
\date{}

\begin{document}
	\maketitle
	
	\hrule
	
	\section*{Problems}
	
	\begin{problem}[3]
		Suppose you have \(m=14\) training examples with \(n=3\) features (excluding the additional all-ones feature for the intercept term, which you should add). The normal equation is \(\theta=(X^TX)^{-1}X^Ty\). For the given values of \(m\) and \(n\), what are the dimensions of \(\theta\), \(X\), and \(y\) in this equation?
	\end{problem}

	\begin{problem}[3]
		The closed form solution for linear regression is \(\theta=(X^TX)^{-1}X^Ty\). Suppose you have \(N=35\) training examples and \(M=5\) features (excluding the bias term). Once the bias term is now included, what are the dimensions of \(X\), \(y\), and \(\theta\) in the closed form equation?
	\end{problem}

	\begin{problem}[2]
		Suppose the scores of the midterm, homework, and final of 4 students are given by the following table.
		
		\begin{table}[ht]
			\centering
			\begin{tabularx}{0.5\textwidth}{|X | X | X|}
				\hline
				\textbf{Midterm} & \textbf{Homework} & \textbf{Final} \\
				\hline
				80 & 85 & 96\\
				72 & 68 & 74\\
				94 & 97 & 90\\
				69 & 70 & 78\\
				\hline
			\end{tabularx}
		\caption{Scores of 4 students in Problem \theproblem.}
		\end{table}
	
	You would like to use linear regression to predict a student's final exam score from his (or her) midterm and homework scores. The normal equation of linear regression is \(\theta=(X^TX)^{-1}X^Ty\). What are \(X\) and \(y\) in this problem?
	\end{problem}

	\begin{problem}[2]
		Given data \(D=\{(x^{(1)},y^{(1)}),\hdots(x^{(N)},y^{(N)})\}\), we obtain \(\hat{w}\), the parameters that minimize the training error cost for the linear regression model \(y=w^Tx\), that we learn from \(D\).
		
		Consider a new dataset \(D_{\text{new}}\) generated by duplicating the points in \(D\) and adding 10 new points that lie along \(y=\hat{w}^Tx\), for a total of \(2N+10\) points. Is the \(\hat{w}_\text{new}\) that we learn for \(y=w^Tx\) from \(D_\text{new}\) equal to \(\hat{w}\)?
	\end{problem}

	\begin{problem}[2]
		Assume that you have a dataset composed of \(N\) observations: the target \(t\), and features \(X\). You want to fit a linear regression model and find the weights \(w\), but you also know that more data is always helpful. Instead of fitting a model with \(t\in\mathbb{R}^n\) and \(X\in\mathbb{R}^{n\times d}\), you concatenate the data and fit a model using \(\displaystyle{\begin{bmatrix}
				t\\
				t
		\end{bmatrix}}\in\mathbb{R}^{2n}\) and \(\displaystyle{\begin{bmatrix}
		X\\
		X
	\end{bmatrix}}\in\mathbb{R}^{2n\times d}\). Will running linear regression on this new dataset give the same weights as on the original dataset?
	\end{problem}

	\begin{problem}[4]
		First, we use a linear regression method to model a set of data. To test our linear regressor, we choose at random some data records to be a training set, and choose at random some of the remaining records to be a test set.
		
		Now let us increase the training size gradually. As the training set increases, what you expect will happen with the mean training and mean testing errors?
	\end{problem}

	\begin{problem}[4]
		We use a multiple linear regression model, which is of the form
		
		\begin{equation*}
			\hat{y}=x\cdot\theta=\sum_{j=0}^{p}\theta_jx_j
		\end{equation*}
	
		Assuming, \(\mathbb{Y}\) is a vector containing our observed responses (\textit{i.e.}, true \(y\)'s); \(\mathbb{X}\) is a design matrix whose first column is 1 (\textit{i.e.}, \(x_0=1\) for all observations), and \(\mathbb{X}_i\) represents the \(i\)-th row of \(\mathbb{X}\). Finally we determing \(\hat{\theta}\) by minimizing average squared loss. \(\hat{\theta}\) is the minimizer of which of the following quantities? (\textit{Select all that apply}.)
		
		\begin{multicols}{4}
			\begin{enumerate}[label=(\alph*)]
				\item \(\displaystyle\sum_{i=1}^n(y_i-\theta)^2\)
				\item \(\dfrac{1}{n}\displaystyle\sum_{i=1}^n(y_i-\theta)^2\)
				\item \(\displaystyle\sum_{i=1}^n(y_i-\mathbb{X}_i\cdot\theta)^2\)
				\item \(\displaystyle{\frac1n||\mathbb{Y}-\mathbb{X}\theta||^2_2}\)
				\item \(\displaystyle{||\mathbb{Y}-\mathbb{X}\theta||^2_2}\)
				\item \(\displaystyle\frac1n\sum_{i=1}^n(y_i-\mathbb{X}_i\cdot\theta)\)
				\item \(\displaystyle\sum_{i=1}^n(y_i-\mathbb{X}_i\cdot\theta)\)
				\item None of the above
			\end{enumerate}
		\end{multicols}
	\end{problem}

	\begin{problem}
		Alice has collected a dataset of dependent and independent variables \(\{(x^{(1)},y^{(1)}),\hdots(x^{(n)},y^{(n)})\}\). She does linear regression on it and obtains the weight vector \(w_\text{Alice}\). Bob also collects the same dataset, but while recording the independent variables (the \(x^{(i)}\)'s) he uses different units. Specifically, Bob's dataset is \(\{(z^{(1)},y^{(1)}),\hdots(z^{(n)},y^{(n)})\}\), where for each \(i\), \(z^{(i)}=cx^{(i)}\ \forall\ c>0\). Bob does linear regression on this dataset and obtains a weight vector \(w_\text{Bob}\).
		
		\begin{enumerate}[label=(\alph*)]
			\item (4 Points) Do Alice and Bob have the same training loss? Is \(w_\text{Alice}=w_\text{Bob}\)? In either case, justify your answer.
			
			\item (6 Points) Suppose now that Bob records each feature in a different unit; \textit{i.e.}, for each coordinate \(j\), \(z_j^{(i)}=c_jx_j^{(i)}\ \forall\ i\in[1,n],c_j>0\). \(c_j\) is a scalar, and those are not all equal. Do Alice and Bob still have the same traning loss and is \(w_\text{Alice}=w_\text{Bob}\)? If yes, justify your answer. If no, provide an example of a dataset where this is not the case.
		\end{enumerate}
	\end{problem}

	\begin{problem}
		Suppose we have a dataset of 100 points whose first few rows are shown below, and that we would like to predict \(\vec{y}\) from \(\vec{v}\) and \(\vec{w}\). Suppose we create a design matrix \(\mathbb{X}\), whose first column is \(\vec{v}\), second column is \(\vec{w}\), and third column \(\vec{u}\) is a new feature \(u_i=|v_i|\). The resulting model is \(\hat{y}_i=\beta_1v_i+\beta_2w_i+\beta_3|v_i|\).
		
		Now consider the following table.
		
		\begin{table}[ht]
			\centering
			\begin{tabularx}{0.5\textwidth}{|X|X|X|}
				\hline
				\(y\)&\(v\)&\(w\)\\
				\hline
				4&\(-30\)&1\\
				6&\(-40\)&2\\
				5&20&3\\
				\hline
			\end{tabularx}
			\caption{Dataset for Problem \theproblem.}
		\end{table}
	
		\begin{enumerate}[label=(\alph*)]
			\item (2 Points) For the data above, suppose we arbitrarily pick \(\vec{\beta}=[0.1,12,0.2]^T\). What is \(\hat{y}_1\)?
			\item (2 Points) For the data above, let \(\vec{e}\) be the residual vector if \(\vec{\beta}=[0.1,12,0.2]^T\). What is \(|e_1|\)?
			\item (2 Points) For the data above, suppose that \(\vec{e}\cdot\vec{e}=9\). What is the mean standard error (MSE)?
			\item (2 Points) For the data above, the matrix \(\mathbb{X}\) has full rank (\textit{i.e.}. no columns are linear combinations of any others). Suppose we compute \(Z=(\mathbb{X}^T\mathbb{X})^{-1}\mathbb{X}^T\vec{y}\). What is \(Z\)? Select one. (\textit{Also fill in the blank(s) where needed}.)
			
			\begin{enumerate}[label=\Roman*.]
				\item It is a vector of length \rule{4em}{0.1px}.
				\item It is a matrix with \rule{4em}{0.1px} rows and \rule{4em}{0.1px} columns.
				\item It does not exist because \(|v_i|\) is not differentiable.
			\end{enumerate}
		\end{enumerate}
	\end{problem}

	\begin{problem}[2]
		State whether the following statements are true or false.
		
		\begin{enumerate}[label=(\alph*),]
			\item In the least-squares linear regression setting, if we double the data matrix \(X\), we double the resulting least squares solution \(\hat{w}\).
			\item When running linear regressions, it is a good idea to look at the largest (in absolute value) regression weights to see which features are most influential in determining the predictions.
		\end{enumerate}
	\end{problem}

	\begin{problem}[2]
		Let \(X\in\mathbb{R}^{m\times n},w\in\mathbb{R}^m\). Consider mean squared error \(L(w)=||Xw-Y||^2_2\). What is \(\nabla_wL(w)\), and what is it used for?
	\end{problem}
	
\end{document}
